\section{Introduction}

Centralized school choice systems have become a defining feature of modern education policy, with nearly 40\% of countries relying on centralized clearinghouses to assign students to secondary schools~\cite{ccas_2024}. These systems are designed to reduce discriminatory admissions practices, lower application costs, improve information provision, and replace fragmented school-level procedures with a transparent and standardized mechanism. By eliminating uncoordinated offers, early contracting, and inefficient reassignment of seats, centralized assignment promises both greater fairness and improved allocative efficiency~\citep{Roth1990}. 

Despite these advantages, the adoption of centralized assignment has been politically contentious. In Chile, for example, opponents argued that the centralized system would undermine parental autonomy and restrict families' ability to choose their children's schools~\citep{rebolledo2018}. Similar debates have emerged in other countries, where concerns about bureaucratic overreach, loss of control, or reduced flexibility have shaped public opinion~\citep{fox5ny_lottery_2022, sfstandard_lottery_2026, guardian_brighton_lottery_2017}. Beyond these political dynamics, the expected effects of centralized school choice on downstream outcomes such as college applications, admissions, and enrollment are unclear. Eliminating selective screening and allowing students to apply through a centralized platform that increases competition at the school level may improve academic preparation~\citep{camposkearns2024} and information provision~\citep{arteaga2022smart}, but it may also weaken school--student match quality~\citep{KaufmanRosenbaum1992, cullen2006, LightStrayer2000}. 

% A growing number of education systems have adopted centralized school assignment mechanisms, often implemented through clearinghouses that coordinate admissions across schools. These systems aim to reduce discriminatory admissions practices, lower application costs, and replace fragmented school-level procedures with transparent and standardized allocation mechanisms \citep{ccas_2024}. By eliminating uncoordinated offers and early contracting, centralized assignment can improve both fairness and allocative efficiency. At the same time, the impact of centralized school choice on students' educational outcomes remains theoretically ambiguous. Removing school-level screening may expand access to higher-quality schools for disadvantaged students, but it may also weaken incentives for academic effort or reduce the quality of student--school matches \citep{RebolledoOlea2018}.

% Empirically identifying these effects is challenging because such reforms are typically implemented once and at scale, often alongside other policy changes, which complicates the construction of credible counterfactuals. As a result, much of the existing literature has focused on short-run outcomes such as sorting patterns, access to preferred schools, or parental satisfaction, leaving open questions about the downstream consequences of centralized assignment for longer-term educational trajectories. \citep{ccas_2024}

Studying the effects of centralized school choice---and the mechanisms through which these reforms operate---is empirically challenging. Reforms of this type are typically implemented once, at a local level, and without natural comparison groups, making it difficult to construct valid counterfactuals and assess their broader implications. Moreover, centralized assignment often coincides with other policy changes---such as funding reforms, accountability shifts, or curricular adjustments---that further complicate identification.  As a result, the existing literature has largely focused on district-level reforms to analyze short-run effects on sorting, access, or satisfaction, leaving open the question of how centralized school choice shapes students' postsecondary trajectories.

In this paper, we address these challenges by exploiting the staggered implementation of Chile's centralized school choice system (SAE), which was rolled out region by region between 2016 and 2020 following the Ministry of Education’s pseudo-random selection of regions for early adoption. 
This staggered rollout provides a unique opportunity to study the causal effect of centralized school choice on college admissions outcomes because it creates quasi-experimental variation in exposure across regions and cohorts. Regions that have not yet implemented the reform provide a natural comparison group for those that have already adopted it, helping disentangle the effect of the policy from broader trends affecting students' higher-education trajectories.
The setting is particularly well suited for this analysis because Chile's college admissions system is itself centralized for the most selective institutions, allowing us to trace how changes in secondary-school assignment propagate into higher-education choices.

Our findings reveal four main results. First, centralized school choice significantly increases the likelihood that students apply to, are admitted to, and enroll in programs within the centralized college admissions system. Second, these effects are heterogeneous: the gains are concentrated among female students, students from high-vulnerability schools, and students from lower-performing schools. Third, the improvements are driven primarily by admissions to less selective and historically under-subscribed programs, with some evidence of displacement in highly selective programs. Finally, the results are robust across alternative specifications, placebo tests, parallel-trends diagnostics, and sample restrictions. Together, these findings show that centralized school choice can meaningfully expand access to higher education, particularly for students historically underserved by decentralized admissions.


The remainder of the paper is organized as follows. 
Section~\ref{sec:literature} discusses the most closely related literature. Section~\ref{sec:background} provides institutional background on Chile's school and college admissions systems. Section~\ref{sec:results} presents the main results. Section~\ref{sec:conclusion} concludes.




